% ============================================
% DOC-24 (ZH)
% 中文翻译 — SPEC-935-R356
% ============================================
\section*{\textbf{DOC-24} --- 保护技术意见书}
\markboth{DOC-24 --- 保护技术意见书}{}

\noindent \textit{从鉴定档案中提取的文件。}

\vspace{0.5em}\noindent\rule{\textwidth}{0.5pt}\vspace{0.5em}


\noindent {\footnotesize\texttt{公共档案馆 $\cdot$ 书证 $\cdot$ 案卷1,260}}\\[-0.2em]

\noindent \textbf{米纳斯吉拉斯州公共档案馆 —— 保护与修复实验室}\\
\noindent \textbf{物品:} 装订笔记本，登记号1,260。\\
\noindent \textbf{负责技师:} 费尔南达·洛佩斯·巴斯托斯，CRT-04/1128。\\
\noindent \textbf{申请事项:} 物理完整性评估与存放建议。\\
\noindent\textbf{日期:} 2026年8月4日。

\noindent \textbf{1. 总体描述}

\noindent 笔记本使用克重不一的纸张（80至120克/平方米），手工鞣制皮革装订，可能为本地或手工制作。总体保存状态与所载日期（1917至1979年间）不一致。纸张没有该时期物品应有的氧化痕迹。载体酸度为中性——对于二十世纪初的非工业化纸张来说并不寻常。

\noindent \textbf{2. 主要发现}

\noindent 在页张分析过程中，发现笔记本厚度与1979年清点登记的页数之间存在差异。原始清点——由安东尼奥·塔瓦雷斯医生在掘尸时编制——列出约190页有文字的页面。当前计数为208张有文字的页张。

\noindent 没有证据表明纸张在物理上被添加。装订完整，没有拆装或重插痕迹。原始缝线，原始胶水。然而，多了18页。

\noindent \textbf{3. 多余页面分析}

\noindent 这18页分布在整本笔记本中，与原始页张交错。每一种情况下，新页——或看似新的页——都插入两页之间，而这两页的文本是连续的，笔迹与紧邻的时期相符。多出的页面仿佛在两者之间“生长”出来，如同笔记本自行生成了新纸，以容纳原始书写中不存在的文字。

\noindent 多出页面的墨迹与原始墨迹化学成分相同，但配方中的铁含量略高——表明氧化时间较晚。相对年代测定测试将这几页的书写时间定位于2020至2026年间。

\noindent 这些多出的页面，写在一本自1979年起就封存于保险柜中的笔记本内部。

\noindent \textbf{4. 多出页面的内容}

\noindent \textit{应管理局要求，18页的内容未被转录。仅告知：笔迹与患者1,260的笔迹不符。多出页面的笔迹初步——未做结论性笔迹分析——与埃托尔·奥利维拉医生最后记录中的笔迹相似。}
\noindent \textbf{5. 技术建议}

\noindent 装入氩气气氛密封袋。每周监测重量与厚度。禁止单人操作。负责技师通过本意见书请求——不再被召集对该物品进行新的分析。

\noindent \textit{费尔南达·洛佩斯·巴斯托斯}\\
\noindent \textit{保护与修复技师}\\
\noindent \textit{CRT-04/1128}

\vspace{0.5em}
\noindent\textit{\textbf{↪ 路线:} \rota{DOC-01} \rotasep \rota{LUC-02} \rotasep \rota{DOC-25}}
